% [claim:C101] [evidence:E009,E012,E013]
\title{A Network-Aware Predictive Digital Twin Framework for Event-Based Production Monitoring}
\titlerunning{Network-Aware Predictive Production Monitoring}
\author{\mbox{}}
\authorrunning{}
\institute{\mbox{}}
\maketitle
\begin{abstract}
Discrete sensor events provide limited evidence for monitoring cyber-physical production systems, especially when observations are delayed, incomplete, or inconsistent with a predictive model. This paper presents a network-aware predictive monitoring framework that separates event prediction from evidence eligibility and anomaly interpretation. A process-specific predictor supplies an expected event-time window, while a reusable monitoring kernel evaluates timestamped binary transitions. Data freshness, edge-history continuity, and calibration applicability determine whether a model-based finding is admissible; unavailable evidence is reported separately from a process mismatch. Physics-based prediction, constant residual correction, and learned residual correction are interchangeable alternatives rather than assumed performance improvements. A conveyor-based MuJoCo cell serves as the single virtual testbed. A frozen Ridge model is assessed with separated training, validation, and calibration data. Its normal-test arrival-time mean absolute error is 15.4 milliseconds versus 18.1 milliseconds for constant correction, but its wider calibration window does not improve fault-warning timeliness. Across 31 recordings and 465 scheduled method replays, both corrected alternatives flag all 18 fault segments under immediate reception, with no findings in nine normal segments. Nine missing-sensor and stalled-drive pairs nevertheless have identical inputs, preventing unique fault attribution. The results support a reusable monitoring formulation and expose limits of learned prediction under incomplete observations. They do not establish transfer across machines, physical diagnostic accuracy, production recovery, or autonomous PLC reconfiguration.
\keywords{Predictive digital twin \and Event-based monitoring \and Data freshness \and Incomplete observations \and Residual learning}
\end{abstract}
